% section: 解けない漸化式と、解くことの意味

\section{解けない漸化式と、解くことの意味}

\subsection{「解けない」は無知という意味ではない}

本書でいう「解ける」とは、主に高校数学で扱う関数と有限回の操作によって一般項を表せることである。
従って「解けない」とは、数学的に何も分からないという意味ではない。

例えば、
\[
  a_{n+1}=a_n^2+1,\qquad a_0=0
\]
なら、
\[
  0,1,2,5,26,677,\ldots
\]
と順番に計算できる。
一般項が簡単な式で書けなくても、次の問いは残っている。

\begin{itemize}
  \item すべての項は正か。
  \item 単調に増加するか。
  \item どの程度の速さで増えるか。
  \item ある範囲に留まるか。
  \item 初期値を変えると、どのように動きが変わるか。
\end{itemize}

一般項は、漸化式を理解するための一つの答えにすぎない。

\subsection{数学では、問いを変えることが解法になる}

一般項が求めにくいとき、同じ問いに固執する必要はない。
\[
  \text{正確な値を求める}
\]
から、
\[
  \text{増加の速さを求める}
\]
へ変えてもよい。
さらに、
\[
  \text{どの対象へ移せば構造が見えるか}
\]
を問うてもよい。

差分は変化を見えるようにした。
ニュートン法は、解そのものではなく解への道を作った。
母関数は無限個の項を一つにまとめた。
行列は複数の数列を一つの状態にまとめた。
チェビシェフ多項式は、三角関数の振動を多項式へ翻訳した。
漸近解析は、正確な値から長期的な傾向へ焦点を移した。

これらは、ばらばらなテクニックではない。
扱いにくい対象を見つけたとき、何を保存し、何を捨て、どの形へ移せばよいかを考えるという、同じ数学的態度である。
